% Chapter 3 outline. The research design, participants, instruments, data,
% algorithms, and statistics below are planning topics, not approved methods.

\newpage
\section{Methodology}
\label{Methodology}
% Provide a reproducible account of how the system and research evaluation are
% conducted. Map every method to a specific objective or research question.

\subsection{Research Design}
% Specify and justify:
% - the BSIT capstone research/development design;
% - the selected software-development process and iteration stages;
% - the study area, partner organization, and deployment context;
% - the simulation/optimization experiment design;
% - the system and user-evaluation designs;
% - participants or expert validators, sampling, inclusion criteria, and sample
%   size when applicable;
% - ethical review, consent, privacy, and operational-safety considerations.
% Required decision: evaluation method and baseline.

\subsection{Data Gathering Tools}
% Describe the tools used to acquire or prepare evidence and system data:
% - formal data requests to the selected DRRM, engineering, barangay, or other
%   official offices;
% - GIS extraction and digitization tools;
% - field or document validation procedures;
% - OpenStreetMap or other external-data acquisition procedures;
% - versioning, metadata, and provenance records.
% QGIS, OpenStreetMap, NAMRIA, and local-government sources are candidates only
% until access, licensing, coverage, and quality are verified.

\subsection{Research Instruments}
% Describe, develop, and validate the instruments actually used, potentially:
% - stakeholder interview or requirements guide;
% - geographic-data quality and completeness checklist;
% - expert scenario and output-validation form;
% - system test cases and expected results;
% - usability or software-quality questionnaire;
% - task-performance or decision-comparison instrument.
% State instrument source/adaptation, scoring, validation, pilot testing, and
% reliability procedures. Do not assume ISO 25010 or another standard unless
% the adviser approves it and the constructs match the research questions.

\subsection{Data Gathering Procedures}
% Present the procedure in chronological, reproducible order:
% 1. Confirm the study area, partner organization, users, and scope.
% 2. Secure permissions and define data-governance responsibilities.
% 3. Inventory and acquire road, map, drainage, hazard, population, center, and
%    resource data required by the approved model.
% 4. Record metadata; clean, georeference, standardize, and validate each layer.
% 5. Define representative scenarios and ground-truth or expert references.
% 6. Design and implement the approved routing, allocation, time, and simulation
%    methods and the user interface.
% 7. Conduct unit, integration, system, scenario, and data-quality tests.
% 8. Conduct expert and/or user evaluation following approved ethics procedures.
% 9. Analyze results and preserve anonymized evidence for reproducibility.

\subsection{Statistical Treatment of Data}
% Select analyses only after the outcome variables and evaluation design are
% approved. Address, as applicable:
% - frequencies, percentages, means/medians, variability, and confidence bounds;
% - route validity, distance/time, capacity violations, allocation feasibility,
%   computation time, and agreement with validated plans;
% - questionnaire scoring, reliability, and interpretation;
% - comparisons with a baseline or alternative method;
% - assumptions, effect sizes, missing data, and decision thresholds.
% Provide formulas, software, significance level, and interpretation rules for
% every statistic actually used. Do not include decorative statistics.
